% ==================================================================
% id:       cond-004
% topic:    הסתברות מותנית ואי-תלות
% subtopic: מאורעות בלתי תלויים
% solution: full
% ==================================================================
זורקים שתי קוביות הוגנות, ונסמן את המאורעות:
$A$ = ״בקובייה הראשונה יצא מספר זוגי״, $B$ = ״בקובייה השנייה יצא מספר זוגי״, $C$ = ״סכום הקוביות זוגי״.
\begin{enumerate}
\item הראו ש־$A,B,C$ בלתי תלויים בזוגות (כלומר, כל זוג מאורעות מתוכם בלתי תלוי).
\begin{solbox}
מרחב המדגם אחיד, $|\Omega|=36$. בכל קובייה יש 3 תוצאות זוגיות מתוך 6, ולכן $\PP(A)=\PP(B)=\frac{18}{36}=\frac12$. סכום הקוביות זוגי כששתיהן זוגיות או ששתיהן אי־זוגיות: $3\cdot3+3\cdot3=18$ תוצאות, ולכן $\PP(C)=\frac12$.

החיתוכים:
\begin{itemize}
\item $A\cap B$ = ״שתי הקוביות זוגיות״: $3\cdot3=9$ תוצאות, ולכן $\PP(A\cap B)=\frac9{36}=\frac14$.
\item $A\cap C$ = ״הראשונה זוגית והסכום זוגי״ = ״שתיהן זוגיות״, ולכן $\PP(A\cap C)=\frac14$.
\item $B\cap C$ = ״השנייה זוגית והסכום זוגי״ = ״שתיהן זוגיות״, ולכן $\PP(B\cap C)=\frac14$.
\end{itemize}
בכל זוג $\frac14=\frac12\cdot\frac12$, כלומר ההסתברות של החיתוך שווה למכפלת ההסתברויות. לכן $A,B,C$ \(\ans{\text{בלתי תלויים בזוגות}}\).
\end{solbox}
\item האם $A,B,C$ בלתי תלויים במשותף?
\begin{solbox}
צריך לבדוק גם את החיתוך של שלושתם. אם שתי הקוביות זוגיות, גם הסכום זוגי, כלומר $A\cap B\subseteq C$, ולכן $A\cap B\cap C=A\cap B$ ו־
\[
\PP(A\cap B\cap C)=\tfrac14\neq\tfrac18=\PP(A)\,\PP(B)\,\PP(C).
\]
לכן הם \(\ans{\text{לא בלתי תלויים במשותף}}\). (במילים: אם ידוע ש־$A$ ו־$B$ קרו, $C$ בוודאות קורה.)
\end{solbox}
\end{enumerate}
